% Generated draft: requires mathematical and source review.
\documentclass[11pt]{article}
\usepackage[margin=25mm]{geometry}
\usepackage{fontspec}
\setmainfont{Latin Modern Roman}
\usepackage{amsmath,amssymb,mathtools,unicode-math}
\setmathfont{Latin Modern Math}
\usepackage{enumitem,needspace,xurl,hyperref}
\setlength{\parindent}{0pt}
\setlength{\parskip}{5pt}
\setlength{\emergencystretch}{3em}
\newcommand{\N}{\mathbb{N}}
\newcommand{\Z}{\mathbb{Z}}
\newcommand{\Q}{\mathbb{Q}}
\newcommand{\R}{\mathbb{R}}
\newcommand{\C}{\mathbb{C}}
\newcommand{\F}{\mathbb{F}}
\newcommand{\A}{\mathbb{A}}
\newcommand{\PP}{\mathbb{P}}
\newcommand{\E}{\mathbb{E}}
\newcommand{\eps}{\varepsilon}
\newcommand{\dd}{\,\mathrm{d}}
\begin{document}
\section*{UnsolvedMath ID: 30003279}
\textbf{Source problem code:} catalogue entry
\subsection{Definitions and mathematical statement}
Let an admissible integer be a positive integer of the form
\[
 n=\prod_{i=1}^{r}p_i,
 \qquad p_i\equiv1\pmod4,
\]
where the \(p_i\) are distinct primes. Allowing the empty product \(n=1\) has no effect on the asymptotic questions below. Set
\[
 R=\sqrt n,
 \qquad L_n=\{(x,y)\in\Z^2:x^2+y^2=n\}.
\]
The source does not define the separation of a configuration. One possible precise interpretation is its Euclidean diameter. Under this interpretation, for an integer \(k\ge2\), define
\[
 D_k(n)=\min_{\substack{A\subseteq L_n\\|A|=k}}
 \max_{u,v\in A}\|u-v\|_2,
\]
with \(D_k(n)=+\infty\) if \(|L_n|<k\). The points in \(A\) are thus distinct.

\textbf{Three-point question, conditional on this interpretation.}
Do there exist constants \(0<c\le C<\infty\) and admissible integers \(n_j\to\infty\), with \(R_j=\sqrt{n_j}\), such that
\[
 cR_j^{1/3}\le D_3(n_j)\le CR_j^{1/3}
 \qquad\text{for every }j?
\]
The source describes \(R^{1/3}\) as a known lower-bound scale. If that description means a uniform lower bound for this particular \(D_3\), the attainability question reduces to whether the displayed upper bound holds along an unbounded sequence. The supplied record does not give the exact lower-bound theorem or its hypotheses.

\textbf{General configuration question.}
For each fixed integer \(k\ge2\), determine sharp lower bounds for \(D_k(n)\) in terms of \(R=\sqrt n\), for admissible \(n\) with \(|L_n|\ge k\), and determine whether matching upper bounds occur along unbounded sequences of admissible integers. Constants may depend on \(k\), but not on the radius along such a sequence.

\textbf{Unresolved qualifications.}
These are explicit candidate formulations, not uniquely determined statements extracted from the record. The source does not specify whether separation means diameter, the length of a containing circular arc, minimum pairwise distance, or another quantity. It also does not specify the intended asymptotic quantifiers or the general \(k\)-point lower bounds. Those choices require mathematical review before either question can be treated as the definitive formulation.
\subsection{Short English statement}
The problem concerns clusters of distinct integer points on circles whose squared radii are products of distinct primes congruent to \(1\) modulo \(4\). It asks whether three-point clusters can attain the stated \(R^{1/3}\) separation scale, and what the sharp scales are for larger clusters. The displayed formulation uses cluster diameter and arbitrarily large radii; the source leaves these choices unresolved.
\subsection{Sources}
\begin{itemize}
\item Local UnsolvedMath registry, supplied complete record for permanent ID \(30003279\), title \emph{Minimal Separations Among Lattice Points on Circles}, exact source code \texttt{catalogue entry}, in the Oberwolfach Reports Open Problems collection. The proposed year is recorded as 2016; no individual proposer is supplied.
\item Supplied report citation: \emph{Analytic Number Theory (2017)}, Oberwolfach Reports, pp.\ 2975--3030, DOI \texttt{10.4171/owr/2016/53}. \url{https://doi.org/10.4171/owr/2016/53}. The citation is reproduced as supplied; its 2017 label and the 2016 report identifier have not been independently reconciled.
\item Supplied literature source, ``OWR 53 report'': \url{https://publications.mfo.de/bitstream/handle/mfo/3557/OWR_2016_53.pdf?isAllowed=y&sequence=1}.
\item Supplied literature source, ``Cilleruelo and Granville'': \url{https://doi.org/10.4153/CJM-2009-057-2}. No title, theorem number, or precise statement is supplied for this reference.
\item The literature assessment dated \textbf{2026-08-21} is supplied dated curation, not independently verified. It labels the problem partially solved and reports substantial results for two, three, and four points, including sharp constructions for four close lattice points, while describing broader close-cluster questions as open. The source explicitly requests independent verification of time-sensitive status and cited claims.
\item The registry labels the statement \texttt{corrected\_verified}, explaining that mathematical content was preserved while notation, typography, or omitted definitions were repaired. This is a supplied editorial label, not independent mathematical verification of this draft.
\end{itemize}
\end{document}
